\begin{scanpage}{0194}{176}
\sectiontitle{28}{Current Waves in an Ideally Conducting Vibrator}
The electrodynamic problem of determining the current in thin cylindrical conductors (vibrator) usually is reduced to an integro-differential equation. The latter is derived by means of boundary conditions on the conductor surface, and is substantially simplified in the case of thin conductors when the inequalities
\[
\frac aL\ll1\quad\text{and}\quad ka\ll1.\tag{28.01}
\]
are fulfilled, where $a$ is the radius and $L$ is the length of the conductor and $k=2\pi/\lambda=\omega/c$.

Its solution may be found, for example, by the method of successive approximations [79, 80] or by the perturbation method [85]. Recently, Vaynshteyn [81, 82] proposed a new solution for this equation. Since we will subsequently base our work on the results of References [81, 82], let us discuss them in more detail.

Let us assume that the vibrator's symmetry axis coincides with the $z$ axis, and its ends have the coordinates $z=z_1$ and $z=z_2$ ($L=z_2-z_1$). In the case of excitation of the dipole by a concentrated external field
\[
E_z^e=\mathcal E\delta(z)\tag{28.02}
\]
the current $J(z)$ in the conductor may obviously be written in the form of the sum of the waves travelling along the conductor with a velocity $c$ from the excitation point $z=0$ and the ends $z=z_1$ and $z=z_2$. In Reference [81] it was shown that the complex amplitudes of these waves are slowly varying functions of the $z$ coordinate. These functions may be approximately expressed in terms of the function $\psi(z)$, so that we obtain the following expression for the current $J(z)$:
\[
J(z)=J_0[\psi(|z|)e^{ik|z|}+A_1\psi(z-z_1)e^{ik(z-z_1)}+A_2\psi(z_2-z)e^{ik(z_2-z)}].\tag{28.03}
\]
\end{scanpage}
